\begin{abstract}
The Fermi Paradox asks why we observe no evidence of extraterrestrial civilizations. We propose that the answer lies in an endogenous mechanism: the dissipation of adaptive potential through intraspecific economic cannibalism---the progressive concentration of resources by elites optimizing extraction over adaptation. We formalize this hypothesis as the EDAP model, a system of nonlinear differential equations governing the co-evolution of technology ($T$), cannibalism ($K$), and cooperation ($C$). The model identifies two distinct attractors: a sustainable development trajectory and a feudal trap---a stable low-technology equilibrium representing civilizational stagnation. We argue that the feudal trap explains the \emph{Great Plateau}: the absence of interstellar expansion. The model further identifies a technofeudal threshold $T_{\text{sing}}$ beyond which automation renders human populations economically redundant to elites, enabling the transition from plateau to \emph{filter}---the complete dissipation of detectable technosignatures. We emphasize that this filter transition is a qualitative implication of the model's structure, not a quantitative prediction; explicit population dynamics and technosignature modeling are required to establish this claim fully. Calibrated via constrained differential evolution on four historical civilizations (Roman Empire, USSR, Russian Federation, United States), the model achieves turn accuracy of 0.70 on training data. For the United States test set, turn accuracy is 0.59, which does not significantly exceed chance (binomial test, $p = 0.31$). We identify structural tension points that align with documented historical interventions by elites and discuss the model's limitations, including a small calibration sample and a trade-off between turn accuracy and early warning lead time.
\end{abstract}
